\section{Exclusions at the counting bound}
\label{sec:exclusions}

We first work one case by hand, then give the data for all five, and
finally describe the exact scan that found them.

\subsection{A complete 31-cap in \texorpdfstring{\(\PG(4,7)\)}{PG(4,7)}}

Suppose \(A\) is a complete \(31\)-cap in \(\PG(4,7)\).  Then \(N=465\),
\(\theta_3=400\), \(\theta_4=2801\), and
\(\Lambda_0=465\cdot6-2801+31=20\); the counting bound is met with a slack
of \(20\), so \(31\) is its smallest admissible value.  The hyperplane
loss polynomial is
\[
 Q(s)=465-400+7\binomtwo{s}-29s=\tfrac12\bigl(7s^2-65s+130\bigr),
\]
with \(Q(0)=65\), \(Q(1)=36\), \(Q(2)=14\), \(Q(3)=-1\), \(Q(6)=-4\),
\(Q(7)=9\), \(Q(8)=29\), and \(Q\) is decreasing before its vertex at
\(s^*=4.64\) and increasing after.  The only values in \([0,20]\) are
\(Q(2)=14\) and \(Q(7)=9\): every solid meets \(A\) in exactly \(2\) or
exactly \(7\) points.

\emph{Global route.}  With \(b_2\) solids of size \(2\) and \(b_7\) of
size \(7\), the first two character equations read
\(b_2+b_7=2801\) and \(2b_2+7b_7=31\cdot400=12400\), so
\(5b_7=6798\), which is not divisible by \(5\).  No such cap exists.

\emph{Local route.}  Fix a secant \(\ell\).  The \(t=\theta_2=57\)
solids through \(\ell\) have \(w_\Pi\in\{0,5\}\) with
\(\sum w_\Pi=R=29\cdot8=232\), and \(232\) is not divisible by \(5\); the
local system is infeasible as well.  Using only the lower half
\(Q(s)\ge0\), the section sizes lie in \(\{0,1,2\}\cup\{7,\dots,31\}\),
so \(w_\Pi\in\{0\}\cup\{5,\dots,29\}\) and the mean \(R/t=4.07\) is
bracketed by \(a=0\) and \(b=5\).  With \(C_0=\binom{29}{2}=406\) and
\(D=7\), Corollary~\ref{cor:closed-form} gives
\(\tau_0=(4\cdot232-812)/14=58/7\), so \(T_\ell\ge9\) on every secant.
Theorem~\ref{thm:secant-local} with \(m=15\) then demands
\(\Lambda_0\ge465\,\varphi_{15}(9)=465\cdot\tfrac9{10}=418.5\), against
\(\Lambda_0=20\).  Either route proves \(t_2(4,7)\ge32\).

\subsection{The five cases}

Table~\ref{tab:cases} gives the same data for all five parameter sets.
In each, \(k\) is the smallest integer satisfying the counting
bound~\eqref{eq:counting-bound}.  The global column reports the
divisibility obstruction in the two-element cases; the local columns
report the one-sided bracket bound and the resulting coverage demand,
which exceeds \(\Lambda_0\) by a wide margin in every case.

\begin{table}[ht]
\centering
\caption{Complete caps excluded at the counting bound.  \(\cS\) is the
admissible set of hyperplane section sizes; the global column gives the
integer that the two-size character equations require to be divisible by
\(s_2-s_1\); \(a,b\) bracket the mean \(R/t\) in the one-sided support
\(\cS_+-2\); \(\beta=\lceil\tau_0\rceil\) is the resulting lower bound on
every \(T_\ell\); the last column is the coverage demand
\(N\varphi_m(\beta)\), to be compared with \(\Lambda_0\).}
\label{tab:cases}
\small
\begin{tabular}{@{}rrrrlrrrrr@{}}
\toprule
\(d\) & \(q\) & \(k\) & \(\Lambda_0\) & \(\cS\) & global & \(a,b\) & \(\tau_0\) & \(\beta\) & \(N\varphi_m(\beta)\)\\
\midrule
4 & 7  & 31  & 20  & \(\{2,7\}\)              & \(6798/5\)     & \(0,5\)   & \(58/7\)      & 9  & \(418.5\)\\
4 & 8  & 37  & 18  & \(\{3,7\}\)              & \(7602/4\)     & \(1,5\)   & \(5/4\)       & 2  & \(444\)\\
4 & 16 & 97  & 32  & \(\{4,9\}\)              & \(144173/5\)   & \(2,7\)   & \(21/4\)      & 6  & \(3990.9\)\\
6 & 8  & 293 & 146 & \(\{29,30,31,43,44\}\)   & ---            & \(29,41\) & \(5009/256\)  & 20 & \(40741.0\)\\
6 & 9  & 387 & 44  & \(\{37,50\}\)            & \(3587183/13\) & \(35,48\) & \(16940/729\) & 24 & \(71703.4\)\\
\bottomrule
\end{tabular}
\end{table}

\begin{proof}[Proof of Theorem~\ref{thm:exclusions}]
For \(\PG(4,7)\) see above.  For \(\PG(4,8)\), \(\PG(4,16)\), and
\(\PG(6,9)\) the admissible set has two elements \(s_1<s_2\), and
\((s_2-s_1)b_{s_2}=k\theta_{d-1}-s_1\theta_d\) is the fraction shown in the
global column, which is not an integer; Proposition~\ref{prop:global-feasibility}
excludes the cap.  For \(\PG(6,8)\) the admissible set is
\(\{29,30,31,43,44\}\), so \(w_\Pi\in\{27,28,29,41,42\}\) and the local
system~\eqref{eq:local-system} is feasible, but the mean
\(R/t=170235/4681=36.37\) is bracketed by \(a=29\) and \(b=41\) with no
admissible value between them;
Corollary~\ref{cor:closed-form} gives \(\tau_0=5009/256\), so every
secant has \(T_\ell\ge20\), and Theorem~\ref{thm:secant-local} with
\(N=42778\), \(m=146\) demands \(\Lambda_0\ge42778\cdot\tfrac{20}{21}>40740\)
against \(\Lambda_0=146\).  In the four two-element cases the one-sided
local route of Table~\ref{tab:cases} gives an independent proof.  The
arithmetic in the table is replayed by the exact-arithmetic script
described below.
\end{proof}

\subsection{The scan and where it stops}

The tests of Sections~\ref{sec:feasibility} and~\ref{sec:coverage} are
finite for every \((d,q,k)\).  A deterministic script, using only exact
integer and rational arithmetic, applies them to every prime power
\(q\le128\) for \(d=3\), \(q\le64\) for \(d=4\), \(q\le16\) for \(d=5\),
and \(q\le9\) for \(d=6\), and to every \(k\) from the counting bound to
twelve above it.  For each \((d,q,k)\) it computes \(\Lambda_0\),
\(\cS\), and the constants~\eqref{eq:local-constants}; decides
feasibility of the first two equations of~\eqref{eq:local-system} by a
dynamic programme over the gaps of \(\cS-2\); solves the global and local
systems exactly, with the congruence, when \(|\cS|\le3\); and evaluates
the bracket bound, the budget test~\eqref{eq:budget}, and the
ceiling~\eqref{eq:ceiling}, in both the two-sided and the one-sided
form.  Its certificate records, for every \((d,q)\), the counting bound,
every excluded \(k\), and the first \(k\) surviving all tests.

The outcome in dimension at least four is exactly
Theorem~\ref{thm:exclusions}: the five parameter sets above are the only
ones in the scanned domain at which the smallest excluded \(k\) equals the
counting bound, and in each of them \(k+1\) survives every test.  (Test
firings above the largest cap of a space, for \(q=2\), are not gains and
are listed separately in the certificate.)  In dimension three the scan
reproduces, for most \(q\), the one-step gain of the secant-local bound
of~\cite[Section~5]{RuddSecantDefects}; in \(\PG(3,7)\) the global
divisibility test excludes \(k=12\), where that bound gives \(13\).  The
limits of the domain in dimensions five and six are set by the running
time of the dynamic programme, not by any mathematical stop.  The script,
its certificate, and the SHA-256 digest of the certificate are part of
the paper's repository; the replay command is in its verification
README.  None of the exclusions depends on the scan: each is a hand
calculation of the kind shown for \(\PG(4,7)\), and the script's role is
to replay them and to establish the negative statement about the rest of
the domain.
